%%%%%%%%%%%%%%%%%%%%%%%%%%%%%%%%%%%%%%%%%%%%%%%%%%%%%%%%%%%%%%%%%%%%%%
% Lớp 10 | Đại số tổ hợp | Nhị thức Newton
% Dạng 0D8?3-2: Khai triển một nhị thức Newton
% Số câu: 34  (sinh tự động bởi tach_cau.py; từ nay thầy sửa trực tiếp file này)
%%%%%%%%%%%%%%%%%%%%%%%%%%%%%%%%%%%%%%%%%%%%%%%%%%%%%%%%%%%%%%%%%%%%%%

%%%%%%%%%%%%%%%% TRẮC NGHIỆM 4 PHƯƠNG ÁN %%%%%%%%%%%%%%%%

% nguồn: MaHoaNH_TracNghiem_Lop10
\begin{ex}%[0D8N3-2]%[ma:e5da4a09]
	Đa thức $P(x)=32 x^5-80 x^4+80 x^3-40 x^2+10 x-1$  là khai triển của nhị thức nào?
	 \choice
		{$(1-2  x)^5$}
		{$(1+2 x)^5$} 
		{\True $(2  x-1)^5$}
		 {$(x-1)^5$}
	\loigiai{ 
		Vì hệ số của $x^5$ là $32$ và dấu trong khai triển đan xen nên  chọn phương án $(2 x-1)^5$.
	} 
\end{ex}

% nguồn: MaHoaNH_TracNghiem_Lop10
\begin{ex}%[0D8N3-2]%[ma:9039c5e8]
	Khai triển của nhị  thức $(x-y)^4$ là
	\choice 
		{\True $x^4-4 x^3 y+6 x^2 y^2-4 x  y^3+y^4$}
		{$x^4-4 x^3 y+4 x^2 y^2-4 x^1  y^3-y^4$}
		{$x^4+4 x^3 y+4  x^2 y^2-4 x^1 y^3+y^4$}
		{$x^4-4 x^3 y-4  x^2 y^2-4 x^1 y^3+y^4$}
	\loigiai{ 
		Ta có $(x-y)^4=x^4-4 x^3  y+6 x^2 y^2-4 x y^3+y^4$.
	 }
\end{ex}

% nguồn: MaHoaNH_TracNghiem_Lop10
\begin{ex}%[0D8N3-2]%[ma:e1410e47]
	Trong khai triển nhị thức Newtơn của $(a+b)^4$, số hạng tổng quát  của khai triển là
	\choice 
		 {$\mathrm{C}_4^{k-1}a^kb^{5-k}$}
		{\True $\mathrm{C}_4^ka^{4-k}b^k$} 
		 {$\mathrm{C}_4^{k+1}a^{5-k}b^{k+1}$}
		{$\mathrm{C}_4^ka^{4-k}b^{4-k}$} 
	\loigiai{ 
		Số  hạng tổng quát của khai triển $(a+b)^4$ là $\mathrm{C}_n^ka^{n-k}b^k=\mathrm{C}_4^ka^{4-k}b^k$.
	} 
\end{ex}

% nguồn: MaHoaNH_TracNghiem_Lop10
\begin{ex}%[0D8N3-2]%[ma:417743ca]
	Khai triển của nhị  thức $(2 x+y)^5$ là
	 \choice
		{$32 x^5+16 x^4 y+8 x^3 y^2+4 x^2 y^3+2  x y^4+y^5$}
		{\True $32  x^5+80 x^4 y+80 x^3 y^2+40 x^2 y^3+10 x y^4+y^5$}
		 {$2 x^5+10 x^4 y+20 x^3 y^2+20 x^2 y^3+10 x y^4+y^5$}
		{$32  x^5+10\,000 x^4 y+80\,000 x^3 y^2+400 x^2 y^3+10 x y^4+y^5$}
	\loigiai{ 
		 Ta có
		\begin{eqnarray*}
			(2  x+y)^5&=&\mathrm{C}_{5}^0(2 x)^5+\mathrm{C}_{5}^1(2 x)^4 y+\mathrm{C}_{5}^2(2 x)^3 y^2+\mathrm{C}_{5}^3(2 x)^2 y^3+\mathrm{C}_{5}^4(2 x) y^4+\mathrm{C}_{5}^5 y^5\\
			&=&32 x^5+80 x^4 y+80 x^3 y^2+40 x^2  y^3+10 x y^4+y^5.
		\end{eqnarray*}
	} 
\end{ex}

% nguồn: MaHoaNH_TracNghiem_Lop10
\begin{ex}%[0D8N3-2]%[ma:2b3d97f6]
	Khai triển  của nhị thức $(x-2)^5$ là
	 \choice
		{$x^5-100 x^4+400 x^3-800  x^2+800 x-32$}
		{$5 x^5-10 x^4+40  x^3-80 x^2+80 x-32$}
		 {\True $x^5-10 x^4+40 x^3-80 x^2+80 x-32$}
		{$x^5+10 x^4+40 x^3+80 x^2+80  x+32$}
	 \loigiai{
		Ta  có
		\begin{eqnarray*}
			(x-2)^5&=&\mathrm{C}_{5}^0 x^5+\mathrm{C}_{5}^1 x^4(-2)+\mathrm{C}_{5}^2 x^3(-2)^2+\mathrm{C}_{5}^3 x^2(-2)^3+\mathrm{C}_{5}^4  x(-2)^4+\mathrm{C}_{5}^5(-2)^5\\
			&=&x^5-10  x^4+40 x^3-80 x^2+80 x-32.
		\end{eqnarray*}
	 }
\end{ex}

% nguồn: MaHoaNH_TracNghiem_Lop10
\begin{ex}%[0D8N3-2]%[ma:af1b7ee3]
	Viết khai  triển nhị thức $(x+1)^5$.
	\choice 
		{\True $x^5 +5x^4 +10x^3  +10x^2 +5x +1$}
		{$x^5 -5x^4 -10x^3 +10x^2 -5x +1$} 
		{$x^5 -5x^4  +10x^3 -10x^2 +5x -1$}
		{$5x^5 +10x^4  +10x^3 +5x^2 +5x +1$}
	\loigiai{ 
		 Ta có
		\begin{eqnarray*}
			(x+1)^5 &=&\mathrm{C}_{5}^0 x^5 +\mathrm{C}_{5}^1 x^4  +\mathrm{C}_{5}^2 x^3 +\mathrm{C}_{5}^3 x^2 +\mathrm{C}_{5}^4 x+\mathrm{C}_{5}^5\\
			&=&x^5+5 x^4+10 x^3+10 x^2+5  x+1.
		\end{eqnarray*}
	} 
\end{ex}

% nguồn: MaHoaNH_TracNghiem_Lop10
\begin{ex}%[0D8N3-2]%[ma:fea86619]
	Khai triển của nhị thức $(1-2  x)^5$ là
	\choice 
		{$5-10 x+40 x^2-80 x^3-80 x^4-32  x^5$}
		{$1+10  x+40 x^2-80 x^3-80 x^4-32 x^5$}
		{\True $1-10 x+40 x^2-80 x^3+80 x^4-32  x^5$}
		{$1+10 x+40 x^2+80  x^3+80 x^4+32 x^5$}
	 \loigiai{
		 Ta có
		\begin{eqnarray*}
			(1-2  x)^5&=&\mathrm{C}_{5}^0+\mathrm{C}_{5}^1(-2 x)^1+\mathrm{C}_{5}^2(-2 x)^2+\mathrm{C}_{5}^3(-2 x)^3+\mathrm{C}_{5}^4(-2 x)^4+\mathrm{C}_{5}^5(-2 x)^5\\
			 &=&1-10 x+40 x^2-80 x^3+80 x^4-32 x^5.
		\end{eqnarray*}
	} 
\end{ex}

% nguồn: MaHoaNH_TracNghiem_Lop10
\begin{ex}%[0D8N3-2]%[ma:6a845949]
	Khai triển của nhị thức  $(3 x+4)^5$ là
	 \choice
		{$x^5+1\,620 x^4+4\,320  x^3+5\,760 x^2+3\,840 x+1\,024$}
		{$243 x^5+405 x^4+4\,320 x^3+5\,760 x^2+3\,840 x+1\,024$} 
		{$243 x^5-1\,620 x^4+4\,320  x^3-5\,760 x^2+3\,840 x-1\,024$}
		 {\True $243 x^5+1\,620 x^4+4\,320 x^3+5\,760 x^2+3\,840 x+1\,024$}
	\loigiai{ 
		Ta  có
		\begin{eqnarray*}
			(3 x+4)^5&=&\mathrm{C}_{5}^0(3 x)^5+\mathrm{C}_{5}^1(3 x)^4 \cdot 4+\mathrm{C}_{5}^2(3 x)^3 \cdot 4^2+\mathrm{C}_{5}^3(3 x)^2 \cdot 4^3+\mathrm{C}_{5}^4(3 x)^1 \cdot 4^4+\mathrm{C}_{5}^5 \cdot  4^5\\
			 &=&243 x^5+1\,620 x^4+4\,320 x^3+5\,760 x^2+3\,840 x+1\,024.
		\end{eqnarray*}
	 }
\end{ex}

% nguồn: MaHoaNH_TracNghiem_Lop10
\begin{ex}%[0D8N3-2]%[ma:08f673f9]
	Khai triển của nhị thức  $(x y+2)^5$ là
	 \choice
		{\True $x^5 y^5+10 x^4 y^4+40 x^3  y^3+80 x^2 y^2+80 x y+32$}
		{$5  x^5 y^5+10 x^4 y^4+40 x^3 y^3+80 x^2 y^2+80 x y+32$}
		{$x^5 y^5+100 x^4  y^4+400 x^3 y^3+80 x^2 y^2+80 x y+32$}
		{$x^5 y^5-10 x^4 y^4+40 x^3 y^3-80 x^2 y^2+80 x  y-32$}
	 \loigiai{
		 Ta có
		\begin{eqnarray*}
			(x y+2)^5&=&\mathrm{C}_{5}^0(x y)^5+\mathrm{C}_{5}^1(x y)^4 \cdot 2^1+\mathrm{C}_{5}^2(x y)^3 \cdot 2^2+\mathrm{C}_{5}^3(x y)^2 \cdot  2^3+\mathrm{C}_{5}^4(x y)^1 \cdot 2^4+\mathrm{C}_{5}^5 \cdot 2^5\\
			&=&x^5 y^5+10 x^4 y^4+40  x^3 y^3+80 x^2 y^2+80 x y+32.
		\end{eqnarray*}
	} 
\end{ex}

% nguồn: MaHoaNH_TracNghiem_Lop10
\begin{ex}%[0D8N3-2]%[ma:67e7ddaf]
	Viết khai triển nhị  thức $(x-y)^5$.
	 \choice
		{\True $x^5-5 x^4 y+10 x^3 y^2-10 x^2 y^3+5  x y^4-y^5$}
		{$x^5+5 x^4 y+10  x^3 y^2+10 x^2 y^3+5 x y^4+y^5$}
		{$x^5-5 x^4 y-10 x^3 y^2-10  x^2 y^3-5 x y^4+y^5$}
		{$x^5+5 x^4  y-10 x^3 y^2+10 x^2 y^3-5 x y^4+y^5$}
	 \loigiai{
		Ta có 
		\begin{eqnarray*}
			(x-y)^5&=&\mathrm{C}_{5}^0 x^5+\mathrm{C}_{5}^1  x^4(-y)+\mathrm{C}_{5}^2 x^3(-y)^2+\mathrm{C}_{5}^3 x^2(-y)^3+\mathrm{C}_{5}^4 x(-y)^4+\mathrm{C}_{5}^5(-y)^5\\
			 &=&x^5-5 x^4 y+10 x^3 y^2-10 x^2 y^3+5 x y^4-y^5.
		\end{eqnarray*}
	} 
\end{ex}

% nguồn: MaHoaNH_TracNghiem_Lop10
\begin{ex}%[0D8H3-2]%[ma:ad23b285]
	Viết khai triển theo công  thức nhị thức Newtơn $(x-1)^5$.
	\choice 
		{$x^5+5x^4+10x^3+10x^2+5x+1$} 
		{$x^5-5x^4-10x^3+10x^2-5x+1$} 
		 {\True $x^5-5x^4+10x^3-10x^2+5x-1$}
		 {$5x^5+10x^4+10x^3+5x^2+5x+1$}
	\loigiai{ 
		 $\left(x-1\right)^5=\mathrm{C}_5^0x^5+\mathrm{C}_5^1x^4+\mathrm{C}_5^2x^3+\mathrm{C}_5^3x^2+\mathrm{C}_5^4x+\mathrm{C}_5^5=x^5+5x^4+10x^3+10x^2+5x+1$.
	} 
\end{ex}

% nguồn: MaHoaNH_TracNghiem_Lop10
\begin{ex}%[0D8H3-2]%[ma:b59a157e]
	Trong khai triển nhị thức Newtơn của  $(2x-3)^4$, số hạng tổng quát của khai triển là
	\choice 
		 {$\mathrm{C}_4^k2^k3^{4-k}\cdot x^{4-k}$}
		{\True $\mathrm{C}_4^k2^{4-k}(-3)^k\cdot  x^{4-k}$}
		{$\mathrm{C}_4^k2^{4-k}3^k\cdot x^{4-k}$} 
		{$\mathrm{C}_4^k2^k\left(-3\right)^{4-k}\cdot  x^{4-k}$}
	 \loigiai{
		Số hạng tổng quát của khai  triển $(2x-3)^4$ là $\mathrm{C}_4^k(2y)^{4-k}\left(-3\right)^k=\mathrm{C}_4^k2^{4-k}(-3)^k\cdot x^{4-k}$.
	 }
\end{ex}

% nguồn: MaHoaNH_TracNghiem_Lop10
\begin{ex}%[0D8H3-2]%[ma:7e8f96e3]
	Khai triển nào sau đây  đúng?
	\choice 
		 {\True $(x+1)^5=x^5+5x^4+10x^3+10x^2+5x+1$}
		{$(x+1)^5=x^5+x^4+x^3+x^2+x+1$} 
		{$(x+1)^5=x^5+5x^4+12x^3+12x^2+5x+1$} 
		 {$(x+1)^5=x^5+4x^4+9x^3+9x^2+4x+1$}
	 \loigiai{
		Ta  có $(x+1)^5=x^5+5x^4+10x^3+10x^2+5x+1$.
	 }
\end{ex}

% nguồn: MaHoaNH_TracNghiem_Lop10
\begin{ex}%[0D8H3-2]%[ma:6f881c03]
	Khai triển $(a + 2)^4$ ta  được kết quả đúng là?
	 \choice
		{\True $\mathrm{C}_{4}^{0}a^4 + \mathrm{C}_{4}^{1}a^32 + \mathrm{C}_{4}^{2}a^22^2 + \mathrm{C}_{4}^{3}a 2^3 +  \mathrm{C}_{4}^{4}2^4$}
		{$\mathrm{C}_{4}^{0}2^4  - \mathrm{C}_{4}^{3}a^12^3 + \mathrm{C}_{4}^{2}a^22^2 - \mathrm{C}_{4}^{3}a^3 2 + \mathrm{C}_{4}^{4}a^4$}
		 {$\mathrm{C}_{4}^{1}a^32 + \mathrm{C}_{4}^{2}a^22^2 + \mathrm{C}_{4}^{3}a2^3 + \mathrm{C}_{4}^{4}a^4$}
		{$\mathrm{C}_{4}^{0}a^4 +  \mathrm{C}_{4}^{3}a^32 - \mathrm{C}_{4}^{2}a^22^2 + \mathrm{C}_{4}^{3}a 2^3 - \mathrm{C}_{4}^{4}2^4$}
	 \loigiai{
		Ta có $(a + 2)^4 = \mathrm{C}_{4}^{0}a^4  + \mathrm{C}_{4}^{1}a^32 + \mathrm{C}_{4}^{2}a^22^2 + \mathrm{C}_{4}^{3}a 2^3 + \mathrm{C}_{4}^{4}2^4$.
	} 
\end{ex}

% nguồn: MaHoaNH_TracNghiem_Lop10
\begin{ex}%[0D8H3-2]%[ma:8a2c52a8]
	Kết  quả của khai triển $(x-y)^4$ là
	 \choice
		{\True $x^4-4x^3y+6x^2y^2-4xy^3+y^4$} 
		{$x^4-x^3y+x^2y^2-xy^3+y^4$} 
		 {$x^4-4x^3y+12x^2y^2-24xy^3+24y^4$}
		 {$x^4+4x^3y+6x^2y^2+4xy^3+y^4$}
	 \loigiai{
		Dựa vào tam giác Pascal hoặc công thức khai triển nhị  thức Niu-tơn ta có ngay
		 \[(x-y)^4=x^4-4x^3y+6x^2y^2-4xy^3+y^4.\]
	 }
\end{ex}

% nguồn: MaHoaNH_TracNghiem_Lop10
\begin{ex}%[0D8H3-2]%[ma:0e260e3f]
	Viết khai triển theo công thức nhị thức  Newtơn $(x-y)^5$.
	 \choice
		{\True  $x^5-5x^4y+10x^3y^2-10x^2y^3+5xy^4-y^5$}
		{$x^5+5x^4y+10x^3y^2+10x^2y^3+5xy^4+y^5$} 
		{$x^5-5x^4y-10x^3y^2-10x^2y^3-5xy^4+y^5$} 
		 {$x^5+5x^4y-10x^3y^2+10x^2y^3-5xy^4+y^5$}
	 \loigiai{
		$\begin{aligned}
		\text{Ta có  }
		 (x-y)^5&=\mathrm{C}_5^0x^5+\mathrm{C}_5^1x^4(-y)+\mathrm{C}_5^2x^3(-y)^2+\mathrm{C}_5^3x^2(-y)^3+\mathrm{C}_5^4x(-y)^4+\mathrm{C}_5^5(-y)^5 \\
		 &=x^5-5x^4y+10x^3y^2-10x^2y^3+5xy^4-y^5.
	\end{aligned}$
} 
\end{ex}

% nguồn: MaHoaNH_TracNghiem_Lop10
\begin{ex}%[0D8H3-2]%[ma:4a4798f7]
	Khai triển của nhị thức $(1-2x)^5$ là 
	\choice 
		 {$5-10x+40x^2-80x^3-80x^4-32x^5$}
		 {$1+10x+40x^2-80x^3-80x^4-32x^5$}
		{\True  $1-10x+40x^2-80x^3+80x^4-32x^5$}
		 {$1+10x+40x^2+80x^3+80x^4+32x^5$}
	 \loigiai{
		$\begin{aligned}
		 \text{Ta có }
		(1-2x)^5&=\mathrm{C}_5^0+\mathrm{C}_5^1(-2x)^1+\mathrm{C}_5^2(-2x)^2+\mathrm{C}_5^3(-2x)^3+\mathrm{C}_5^4(-2x)^4+\mathrm{C}_5^5(-2x)^5 \\ 
		 &=1-10x+40x^2-80x^3+80x^4-32x^5.
	\end{aligned}$
 }
\end{ex}

% nguồn: MaHoaNH_TracNghiem_Lop10
\begin{ex}%[0D8H3-2]%[ma:8d210065]
	Khai triển của nhị thức $(3x+4)^5$ là 
	 \choice
		 {$x^5+1620x^4+4320x^3+5760x^2+3840x+1024$}
		{$243x^5+405x^4+4320x^3+5760x^2+3840x+1024$} 
		{$243x^5-1620x^4+4320x^3-5760x^2+3840x-1024$} 
		 {\True $243x^5+1620x^4+4320x^3+5760x^2+3840x+1024$}
	 \loigiai{
		$\begin{aligned}
		 \text{Ta có }
		(3x+4)^5&=\mathrm{C}_5^0(3x)^5+\mathrm{C}_5^1(3x)^4\cdot 4+\mathrm{C}_5^2(3x)^3\cdot 4^2+\mathrm{C}_5^3(3x)^2\cdot 4^3 +\mathrm{C}_5^4(3x)^1\cdot 4^4  +\mathrm{C}_5^5\cdot 4^5 \\
		&=243x^5+1620x^4+4320x^3+5760x^2+3840x+1024. 
	\end{aligned}$
} 
\end{ex}

% nguồn: MaHoaNH_TracNghiem_Lop10
\begin{ex}%[0D8H3-2]%[ma:f558feaa]
	Khai triển của nhị thức $\left(x-\dfrac{1}{x}\right)^5$  là
	 \choice
		{$x^5+5 x^3+10 x+\dfrac{10}{x}+\dfrac{5}{x^3}+\dfrac{1}{x^5}$} 
		{\True $x^5-5 x^3+10  x-\dfrac{10}{x}+\dfrac{5}{x^3}-\dfrac{1}{x^5}$}
		{$5  x^5-10 x^3+10 x-\dfrac{10}{x}+\dfrac{5}{x^3}-\dfrac{1}{x^5}$}
		 {$5 x^5+10 x^3+10 x+\dfrac{10}{x}+\dfrac{5}{x^3}+\dfrac{1}{x^5}$}
	\loigiai{ 
		Ta  có
		\begin{eqnarray*}
			\left(x-\dfrac{1}{x}\right)^5&=&\mathrm{C}_{5}^0  x^5+\mathrm{C}_{5}^1 x^4 \left(\dfrac{-1}{x}\right)^1+\mathrm{C}_{5}^2 x^3\left(\dfrac{-1}{x}\right)^2+\mathrm{C}_{5}^3 x^2\left(\dfrac{-1}{x}\right)^3+\mathrm{C}_{5}^4 x^1\left(\dfrac{-1}{x}\right)^4+\mathrm{C}_{5}^5\left(\dfrac{-1}{x}\right)^5\\
			&=&x^5-5  x^3+10 x-\dfrac{10}{x}+\dfrac{5}{x^3}-\dfrac{1}{x^5}.
		\end{eqnarray*}
	} 
\end{ex}

% nguồn: MaHoaNH_TracNghiem_Lop10
\begin{ex}%[0D8H3-2]%[ma:114f101a]
	 Khai triển của nhị thức $(x-2)^5$.
	\choice 
		 {$x^5-100x^4+400x^3-800x^2+800x-32$}
		{$5x^5-10x^4+40x^3-80x^2+80x-32$} 
		{\True  $x^5-10x^4+40x^3-80x^2+80x-32$}
		{$x^5+10x^4+40x^3+80x^2+80x+32$} 
	 \loigiai{
		$\begin{aligned}
		\text{Ta có  }
		(x-2)^5&=\mathrm{C}_5^0x^5+\mathrm{C}_5^1x^4(-2)+\mathrm{C}_5^2x^3(-2)^2+\mathrm{C}_5^3x^2(-2)^3+\mathrm{C}_5^4x(-2)^4+\mathrm{C}_5^5(-2)^5  \\
		&=x^5-10x^4+40x^3-80x^2+80x-32. 
	\end{aligned}$
 }
\end{ex}

% nguồn: MaHoaNH_TracNghiem_Lop10
\begin{ex}%[0D8H3-2]%[ma:a75a0985]
	Đa thức $P(x)=x^5-5 x^4  y+10 x^3 y^2-10 x^2 y^3+5 x y^4-y^5$ là khai triển của nhị thức nào dưới đây?
	\choice 
		{$(2x-y)^5$} 
		{$(x-2 y)^5$} 
		{$(x+y)^5$} 
		 {\True $(x-y)^5$}
	\loigiai{ 
		 Ta có
		\begin{eqnarray*}
			 (x-y)^5&=&\mathrm{C}^0_5x^5+\mathrm{C}^1_5x^4(-y)+\mathrm{C}^2_5x^3(-y)^2+\mathrm{C}^3_5x^2(-y)^3+\mathrm{C}^4_5x(-y)^4+\mathrm{C}^5_5(-y)^5\\
			&=&x^5-5 x^4 y+10 x^3 y^2-10 x^2 y^3+5 x y^4-y^5. 
		\end{eqnarray*}
	} 
\end{ex}

% nguồn: MaHoaNH_TracNghiem_Lop10
\begin{ex}%[0D8H3-2]%[ma:cec5d0a3]
	Viết khai triển theo công thức  nhị thức Niu-tơn $\left(x-y\right)^5$ là
	 \choice
		 {$x^5+5x^4y+10x^3y^2+10x^2y^3+5xy^4+y^5$}
		 {$x^5-5x^4y-10x^3y^2-10x^2y^3-5xy^4+y^5$}
		{\True  $x^5-5x^4y+10x^3y^2-10x^2y^3+5xy^4-y^5$}
		{$x^5+5x^4y-10x^3y^2+10x^2y^3-5xy^4+y^5$} 
	\loigiai{ 
		 Ta có
		\begin{eqnarray*}
			 &\left(x-y\right)^5&=\mathrm{C}_{5}^0x^5\left(-y\right)^0+\mathrm{C}_{5}^1x^4\left(-y\right)^1+\mathrm{C}_{5}^2x^3\left(-y\right)^2+\mathrm{C}_{5}^3x^2\left(-y\right)^3+\mathrm{C}_{5}^4x^1\left(-y\right)^4+\mathrm{C}_{5}^5x^0\left(-y\right)^5\\
			 &&=x^5-5x^4y+10x^3y^2-10x^2y^3+5xy^4-y^5.
		\end{eqnarray*}
	 }
\end{ex}


%%%%%%%%%%%%%%%% ĐÚNG/SAI %%%%%%%%%%%%%%%%

% nguồn: MaHoaNH_DungSai_Lop10
\begin{ex}%[0D8H3-2]%[ma:73f13474]
	Biết khai triển nhị thức $(x-2)^4=ax^4+bx^3+cx^2+dx+e$ với  $a,b,c,d,e$ là các số thực. Hệ thức nào sau đâu là đúng?
	 \choiceTF
		{\True $a+c=25$} 
		{$a+e=16$} 
		{\True  $a+b+c+d+e=1$}
		 {$c+d=8$}
	\loigiai{ 
		Ta có $\displaystyle  (x-2)^4 =\sum \limits_{k=0}^4 \mathrm{C}_4^k \cdot x^{4-k} \cdot (-2)^k =x^4-8x^3+24x^2-32x+16$. Khi đó, ta có
		\begin{itemize}
			\item $c+d=-8$. 
			\item  $a+e=17$.
			\item $a+c=25$. 
			 \item $a+b+c+d+e=1$.
		\end{itemize}
	} 
\end{ex}

% nguồn: MaHoaNH_DungSai_Lop10
\begin{ex}%[0D8H3-2]%[ma:6f2edad4]
	Trong khai  triển ${\left(xy+\dfrac{1}{y}\right)}^5$ có số hạng nào trong các số hạng sau?
	\choiceTF 
		 {\True $5x^4y^3$}
		{$5x^3y$} 
		 {\True $10x^3y$}
		{$10x^3y$} 
	\loigiai{ 
		Theo  khai triển nhị thức Niu-tơn, ta có\\
		$${\left(xy+\dfrac{1}{y}\right)}^5=\displaystyle \sum\limits_{k=0}^5{\mathrm{C}_5^k\cdot {\left(xy\right)}^{5-k}\cdot {\left(\dfrac{1}{y}\right)}^k}=\displaystyle  \sum\limits_{k=0}^5{\mathrm{C}_5^k\cdot {x}^{5-k}\cdot {y}^{5-2k}}.$$
		Hệ  số của $x^3y$ ứng với $\heva{
		 & 5-k=3 \\
		& 5-2k=1 \\}\Leftrightarrow  k=2\xrightarrow{}$ số hạng cần tìm $\mathrm{C}_5^2x^3y=10x^3y.$\\
		 Hệ số của $x^4y^3$ ứng với $\heva{
		&  5-k=4 \\
		& 5-2k=3 \\}\Leftrightarrow k=1\xrightarrow{}$  số hạng cần tìm $\mathrm{C}_5^1x^4y^3=5x^4y^3.$
	 }
\end{ex}

% nguồn: MaHoaNH_DungSai_Lop10
\begin{ex}%[0D8V3-2]%[ma:dcd45e6f]
	Cho hai biểu  thức $f(x)=(2x+1)^5$ và $g(x)=(x-1)^6$.
	\choiceTF{Khai triển $f(x)$ thành đa thức  có $5$ số hạng}
		{Hệ số của  $x^4$ trong khai triển $g(x)$ thành đa thức là $20$}
		 {Tổng các hệ số trong trong khai triển $g(x)$ thành đa thức bằng $-1$}
		{\True Hệ số của $x^2$ trong khai triển biểu thức $h(x)=f(x)+g(x)$ thành  đa thức là $55$}
	\loigiai{ 
		 $f(x)=(2x+1)^5=32x^5+80x^4+80x^3+40x^2+10x+1$.\\
		 $g(x)=(x-1)^6=x^6-6x^5+15x^4-20x^3+15x^2-6x+1$.
		\begin{itemchoice}
			\itemch Khai  triển $f(x)$ thành đa thức có $6$ số hạng
			\itemch Hệ số của $x^4$ trong khai triển $g(x)$ thành đa thức là  $15$.
			\itemch Tổng các hệ số trong khai triển $g(x)$ thành đa thức bằng $0$. 
			\itemch Ta có $h(x)=x^6+26x^5+95x^4+60x^3+55x^2+4x+2$.\\ 
			Hệ số của $x^2$ trong  khai triển biểu thức $h(x)=f(x)+g(x)$ thành đa thức là $55$.
		\end{itemchoice}
	} 
\end{ex}


%%%%%%%%%%%%%%%% TRẢ LỜI NGẮN %%%%%%%%%%%%%%%%

% nguồn: MaHoaNH_TraLoiNgan_Lop10
\begin{bt}%[0D8C3-2]%[ma:bff9e111]
	Ông A có $500$ triệu đồng và ông B có $600$ triệu đồng gửi hai ngân hàng khác nhau với lãi suất lần lượt là $6\%$/năm và $4\%$/năm. Với giả thiết sau mỗi tháng người đó không rút tiền thì số tiền lãi được nhập vào số tiền ban đầu. Đây được gọi là hình thức lãi kép. Biết số tiền cả vỗn  lẫn lãi $P_n$ sau $n$ tháng được tính bởi công thức $P_n = P_0(1+r)^n$, trong đó $P_0$ là số tiền gửi lúc đầu và $r$ là lãi suất của một tháng. Dùng hai số hạng đầu tiên trong khai triển của nhị thức Newtơn, ước lượng đến tháng $a$ của năm thứ $b$ (tính từ lúc hai ông bắt đầu gửi) thì số tiền thu được của hai ông là bằng nhau. Tính $ab$.
	 \shortans{156}
	 \loigiai{
	 Đổi lãi suất: $6\%$/năm $=0{,}5\%$/tháng và $4\%$/năm $=\dfrac{1}{3}\%$/tháng.\\
	Dùng hai số hạng đầu tiên trong khai  triển của nhị thức Newtơn, ta ước lượng được như sau:
	\begin{itemize}
		\item Số tiền ông A thu được sau $n$  tháng là
		 $$
		500\left(1 + \dfrac{1}{200}\right)^n \approx  500\left(\mathrm{C}^{0}_{n} + \mathrm{C}^{1}_{n}\dfrac{1}{200}\right) = 500 + \dfrac{5n}{2}.
		 $$
		 \item Số tiền ông B thu được sau $n$ tháng là
		 $$
		 600\left(1 + \dfrac{1}{300}\right)^n \approx 600\left(\mathrm{C}^{0}_{n} + \mathrm{C}^{1}_{n}\dfrac{1}{300}\right) = 600 + 2n.
		$$ 
	\end{itemize}
	Để số  tiền của hai ông thu được là bắng nhau thì $500 + \dfrac{5n}{2} = 600 + 2n$ hay $n=200$.\\
	Do $\dfrac{200}{12} =  16\dfrac{8}{12}$ nên đến tháng thứ $8$ của năm thứ $17$ thì số tiền của hai ông sẽ bằng nhau.\\
	 Vậy $ab= 156$.
	 }
\end{bt}


%%%%%%%%%%%%%%%% TỰ LUẬN %%%%%%%%%%%%%%%%

% nguồn: MaHoaNH_TuLuan_Lop10
\begin{bt}%[0D8N3-2]%[ma:59967722]
	 Khai triển các đa thức sau
	\begin{enumEX}{3}
		\item  $(x+3)^{4}$
		 \item $(3x+2y)^{4}$
		\item  $(x+5)^{4}+(x-5)^{4}$
		\item $(x-2y)^{5}$ 
		 \item $(x+y)^{4}$
		\item  $(1+x)^{4}$
		\item $(x-y)^{5}$ 
		\item $(x+1)^{5}$ 
		 \item $(2x-3)^{4}$
		\item $\left( x^{2}+\dfrac{1}{x}\right) ^{4}$ 
	\end{enumEX}
	\loigiai{ 
		\begin{enumerate}[a)]
			\item  $(x+3)^{4} \\
			= \mathrm{C}_{4}^{0}x^{4} +  \mathrm{C}_{4}^{1}x^{3} \cdot 3 + \mathrm{C}_{4}^{2}x^{2} \cdot 3^{2} + \mathrm{C}_{4}^{3}x \cdot 3^{3} + \mathrm{C}_{4}^{4}3^{4} = x^{4}+12x^{3}+54x^{2}+108x+81$.
			\item $(3 x+2 y)^4\\ 
			 = \mathrm{C}_4^0(3 x)^4+\mathrm{C}_4^1(3 x)^3 \cdot 2 y+\mathrm{C}_4^2(3 x)^2 \cdot(2 y)^2+\mathrm{C}_4^3(3 x) \cdot(2 y)^3+\mathrm{C}_4^4(2 y)^4 \\
			= 81 x^4+216 x^3 y+216 x^2 y^2+96 x y^3+16 y^4.$ 
			\item  $(x+5)^{4}+(x-5)^{4}$
			\newline $= (\mathrm{C}_{4}^{0}x^{4} + \mathrm{C}_{4}^{1}x^{3}  \cdot 5 + \mathrm{C}_{4}^{2}x^{2} \cdot 5^{2} + \mathrm{C}_{4}^{3}x \cdot 5^{3} + \mathrm{C}_{4}^{4} \cdot 5^{4}) + (\mathrm{C}_{4}^{0}x^{4} - \mathrm{C}_{4}^{1}x^{3} \cdot 5 + \mathrm{C}_{4}^{2}x^{2} \cdot 5^{2} - \mathrm{C}_{4}^{3}x \cdot 5^{3} + \mathrm{C}_{4}^{4} \cdot 5^{4})$
			 \newline $= 2\mathrm{C}_{4}^{0}x^{4} + 2\mathrm{C}_{4}^{2}x^{2} \cdot 5^{2} + 2\mathrm{C}_{4}^{4} \cdot 5^{4} = 2x^{4}+300x^{2}+1250$.
			 \item $(x-2y)^{5}\\
			= \mathrm{C}_{5}^{0}x^{5} - \mathrm{C}_{5}^{1}x^{4}(2y) +  \mathrm{C}_{5}^{2}x^{3}(2y)^{2} - \mathrm{C}_{5}^{3}x^{2}(2y)^{3} + \mathrm{C}_{5}^{4}x(2y)^{4} - \mathrm{C}_{5}^{5}(2y)^{5}\\
			 =x^{5}-10x^{4}y+40x^{3}y^{2}-80x^{2}y^{3}+80xy^{4}-32y^{5}$.
			\item  $(x+y)^{4}\\
			= \mathrm{C}_{4}^{0}x^{4} + \mathrm{C}_{4}^{1}x^{3}y + \mathrm{C}_{4}^{2}x^{2}y^{2}  + \mathrm{C}_{4}^{3}xy^{3} + \mathrm{C}_{4}^{4}y^{4}$
			\item  $(1+x)^{4}\\
			= \mathrm{C}_{4}^{0}1^{4}  + \mathrm{C}_{4}^{1}1^{3}x + \mathrm{C}_{4}^{2}1^{2}x^{2} + \mathrm{C}_{4}^{3}1x^{3} + \mathrm{C}_{4}^{4}x^{4} = 1+4x+6x^{2}+4x^{3}+x^{4}$.
			\item $(x-y)^{5}\\ 
			=  x^{5}-5x^{4}y + 10x^{3}y^{2}-10x^{2}y^{3}+5xy^{4}-y^{5}$.
			 \item $(x+1)^{5}\\
			=x^{5}+5x^{4}+10x^{3}+10x^{2}+5x+1.$ 
			\item  $(2x-3)^{4}\\
			 = \mathrm{C}_{4}^{0}(2x)^{4} + \mathrm{C}_{4}^{1}(2x)^{3}(-3) + \mathrm{C}_{4}^{2}(2x)^{2}(-3)^{2} + \mathrm{C}_{4}^{3}(2x)(-3)^{3} + \mathrm{C}_{4}^{4}(-3)^{4}\\
			=16x^{4}-96x^{3}+216x^{2}-216x+81$. 
			 \item $\left( x^{2}+\dfrac{1}{x}\right) ^{4}\\
			 = \mathrm{C}_{4}^{0}\left( x^{2}\right) ^{4} + \mathrm{C}_{4}^{1}(x^{2})^{3}\left( \dfrac{1}{x}\right) + \mathrm{C}_{4}^{2}(x^{2})^{2}\left( \dfrac{1}{x}\right) ^{2} + \mathrm{C}_{4}^{3}(x^{2})\left( \dfrac{1}{x}\right) ^{3} + \mathrm{C}_{4}^{4}\left( \dfrac{1}{x}\right)^{4}\\
			= \mathrm{C}_{4}^{0}x^{8} + \mathrm{C}_{4}^{1}x^{6}\dfrac{1}{x} + \mathrm{C}_{4}^{2}x^{4}\dfrac{1}{x^{2}} +  \mathrm{C}_{4}^{3}x^{2}\dfrac{1}{x^{3}} + \mathrm{C}_{4}^{4}\dfrac{1}{x^{4}}\\
			 = x^{8}+4x^{5}+6x^{2}+\dfrac{4}{x} + \dfrac{1}{x^{4}}$.
		\end{enumerate}
	} 
\end{bt}

% nguồn: MaHoaNH_TuLuan_Lop10
\begin{bt}%[0D8H3-2]%[ma:066ab413]
	Biểu diễn $\left( 3+\sqrt{2}\right)  ^{5}-\left( 3-\sqrt{2}\right) ^{5}$ dưới dạng $a+b\sqrt{2}$ với $a$, $b$ là các số nguyên.
	\loigiai{ 
		Ta có $\left( 3+\sqrt{2}\right) ^{5} = \mathrm{C}_{5}^{0}3^{5} + \mathrm{C}_{5}^{1}3^{4}\sqrt{2} + \mathrm{C}_{5}^{2}3^{3}\left( \sqrt{2}\right)  ^{2} + \mathrm{C}_{5}^{3}3^{2}\left( \sqrt{2}\right) ^{3} + \mathrm{C}_{5}^{4}3\left( \sqrt{2}\right) ^{4} + \mathrm{C}_{5}^{5}\left( \sqrt{2}\right) ^{5}$\\
		và  $\left( 3-\sqrt{2}\right) ^{5} = \mathrm{C}_{5}^{0}3^{5} - \mathrm{C}_{5}^{1}3^{4}\sqrt{2} + \mathrm{C}_{5}^{2}3^{3}\left( \sqrt{2}\right) ^{2} - \mathrm{C}_{5}^{3}3^{2}\left( \sqrt{2}\right) ^{3} + \mathrm{C}_{5}^{4}3\left( \sqrt{2}\right) ^{4} - \mathrm{C}_{5}^{5}\left( \sqrt{2}\right) ^{5}$\\
		Suy  ra \\
		$\left( 3+\sqrt{2}\right) ^{5}-\left( 3-\sqrt{2}\right) ^{5}\\ 
		=  2 \left[\mathrm{C}_{5}^{1}3^{4}\sqrt{2} + \mathrm{C}_{5}^{3}3^{2}\left( \sqrt{2}\right) ^{3} + \mathrm{C}_{5}^{5}\left( \sqrt{2}\right) ^{5} \right]\\
		= 2 \left[405\sqrt{2} + 180\sqrt{2} +  4\sqrt{2} \right]= 2 \left[589\sqrt{2} \right] = 1178\sqrt{2}$.\\
		Vậy  $(3+\sqrt{2})^{5}-(3-\sqrt{2})^{5} = 0 + 1178\sqrt{2}$.
	 }
\end{bt}

% nguồn: MaHoaNH_TuLuan_Lop10
\begin{bt}%[0D8H3-2]%[ma:6f57347d]
	Khai triển thành đa thức của $\left( 2-3x\right) ^{2n}$ dưới dạng tổng quát, biết $n$ là  số nguyên dương thỏa mãn $\mathrm{C}_{2n+1}^{1}+\mathrm{C}_{2n+1}^{3}+\mathrm{C}_{2n+1}^{5}+\ldots+\mathrm{C}_{2n+1}^{2n+1}=1024$.
	\loigiai{ 
		Ta có tính chất của tổ hợp $\mathrm{C}_{m}^{1} + \mathrm{C}_{m}^{3} + \mathrm{C}_{m}^{5} + \ldots  = 2^{m-1}$. \\
		Áp dụng với $m = 2n+1$, ta  có
		$\mathrm{C}_{2n+1}^{1}+\mathrm{C}_{2n+1}^{3}+\mathrm{C}_{2n+1}^{5}+\ldots+\mathrm{C}_{2n+1}^{2n+1} = 2^{(2n+1)-1}  = 2^{2n}$.\\
		Theo đề bài  $2^{2n} = 1024$.
		Mà  $1024 = 2^{10}$.
		Suy ra $2n  = 10 \Rightarrow n=5$.\\
		Vậy,  ta cần khai triển $(2-3x)^{2n} = (2-3x)^{10}$.\\
		Theo công  thức nhị thức Newton
		$(2-3x)^{10} = \displaystyle\sum \limits_{k=0}^{10} \mathrm{C}_{10}^{k} 2^{10-k} (-3x)^{k} = \displaystyle\sum \limits_{k=0}^{10} \mathrm{C}_{10}^{k} 2^{10-k}  (-3)^{k} x^{k}$.
	 }
\end{bt}

% nguồn: MaHoaNH_TuLuan_Lop10
\begin{bt}%[0D8H3-2]%[ma:77752c11]
	\begin{enumerate}[a)]
		\item Dùng hai số hạng  đầu tiên trong khai triển của $(1+0{,}05)^{4}$ để tính giá trị gần đúng của $1{,}05^{4}$.
		\item Dùng máy tính cầm tay tính giá trị của $1{,}05^{4}$ và  tính sai số tuyệt đối của giá trị gần đúng nhận được ở câu a.
	\end{enumerate}
	 \loigiai{
		\begin{enumerate}[a)]
			\item Ta có $(1+x)^n = \mathrm{C}_n^0 1^n x^0 + \mathrm{C}_n^1  1^{n-1} x^1 + \mathrm{C}_n^2 1^{n-2} x^2 + \dots + \mathrm{C}_n^n x^n$.\\
			Với $(1+0,05)^{4}$, hai  số hạng đầu tiên là
			$\mathrm{C}_{4}^{0} \cdot 1^4 \cdot  (0,05)^0 + \mathrm{C}_{4}^{1} \cdot 1^3 \cdot (0,05)^1 = 1 + 0,2 = 1{,}2$.\\
			Vậy  $1{,}05^{4} \approx 1{,}2$.
			\item Giá  trị của $1{,}05^{4}$ tính bằng máy tính là $1{,}21550625$.\\
			Sai số tuyệt đối  = $| \text{Giá trị đúng} - \text{Giá trị gần đúng} |$
			$= |1{,}21550625 - 1{,}2|  = 0{,}01550625$.
		\end{enumerate}
	 }
\end{bt}

% nguồn: MaHoaNH_TuLuan_Lop10
\begin{bt}%[0D8H3-2]%[ma:efdc8c9f]
	Khai triển  các đa thức sau:
	\begin{multicols}{2}
		\begin{enumerate}[a)]
			 \item $(x+1)^{5}$
			 \item $(x-1)^{5}$
			 \item $(x+2)^{5}$
			\item  $(2x+y)^{5}$
			 \item $(x-3y)^{5}$
			 \item $(2x+3y)^{5}$
		\end{enumerate}
	\end{multicols}
	\loigiai{ 
		\begin{enumerate}[a)]
			\item $(x+1)^{5}\\ 
			= \mathrm{C}_5^0 x^5  1^0 + \mathrm{C}_5^1 x^4 1^1 + \mathrm{C}_5^2 x^3 1^2 + \mathrm{C}_5^3 x^2 1^3 + \mathrm{C}_5^4 x^1 1^4 + \mathrm{C}_5^5 x^0 1^5\\
			 = x^{5}+5x^{4}+10x^{3}+10x^{2}+5x+1$.
			\item $(x-1)^{5}\\ 
			= \mathrm{C}_5^0 x^5 (-1)^0 + \mathrm{C}_5^1 x^4 (-1)^1 + \mathrm{C}_5^2 x^3 (-1)^2 +  \mathrm{C}_5^3 x^2 (-1)^3 + \mathrm{C}_5^4 x^1 (-1)^4 + \mathrm{C}_5^5 x^0 (-1)^5\\
			= x^{5}-5x^{4}+10x^{3}-10x^{2}+5x-1$. 
			\item  $(x+2)^{5}\\
			= \mathrm{C}_5^0 x^5  2^0 + \mathrm{C}_5^1 x^4 2^1 + \mathrm{C}_5^2 x^3 2^2 + \mathrm{C}_5^3 x^2 2^3 + \mathrm{C}_5^4 x^1 2^4 + \mathrm{C}_5^5 x^0 2^5\\
			= 1 \cdot x^5 \cdot  1 + 5 \cdot x^4 \cdot 2 + 10 \cdot x^3 \cdot 4 + 10 \cdot x^2 \cdot 8 + 5 \cdot x \cdot 16 + 1 \cdot 1 \cdot 32\\
			=  x^{5}+10x^{4}+40x^{3}+80x^{2}+80x+32$.
			\item $(2x+y)^{5}\\ 
			= \mathrm{C}_5^0 (2x)^5 y^0 + \mathrm{C}_5^1 (2x)^4 y^1 + \mathrm{C}_5^2 (2x)^3 y^2 + \mathrm{C}_5^3 (2x)^2 y^3 + \mathrm{C}_5^4 (2x)^1 y^4 + \mathrm{C}_5^5 (2x)^0 y^5\\ 
			= 1 \cdot 32x^5 \cdot 1 + 5 \cdot 16x^4 \cdot y + 10 \cdot 8x^3 \cdot y^2 +  10 \cdot 4x^2 \cdot y^3 + 5 \cdot 2x \cdot y^4 + 1 \cdot 1 \cdot y^5\\
			=  32x^{5}+80x^{4}y+80x^{3}y^{2}+40x^{2}y^{3}+10xy^{4}+y^{5}$.
			\item $(x-3y)^{5}\\ 
			= \mathrm{C}_5^0 x^5 (-3y)^0 + \mathrm{C}_5^1 x^4 (-3y)^1 + \mathrm{C}_5^2 x^3 (-3y)^2 + \mathrm{C}_5^3 x^2  (-3y)^3 + \mathrm{C}_5^4 x^1 (-3y)^4 + \mathrm{C}_5^5 x^0 (-3y)^5\\
			 = 1 \cdot x^5 \cdot 1 + 5 \cdot x^4 \cdot (-3y) + 10 \cdot x^3 \cdot (9y^2) + 10 \cdot x^2 \cdot (-27y^3) + 5 \cdot x \cdot (81y^4) + 1 \cdot 1 \cdot (-243y^5)\\
			 = x^{5}-15x^{4}y+90x^{3}y^{2}-270x^{2}y^{3}+405xy^{4}-243y^{5}$.
			\item $(2x+3y)^{5}\\ 
			= \mathrm{C}_5^0 (2x)^5 (3y)^0 + \mathrm{C}_5^1 (2x)^4 (3y)^1 + \mathrm{C}_5^2 (2x)^3 (3y)^2 + \mathrm{C}_5^3 (2x)^2 (3y)^3 + \mathrm{C}_5^4  (2x)^1 (3y)^4 + \mathrm{C}_5^5 (2x)^0 (3y)^5\\= 1 \cdot 32x^5 \cdot 1 + 5 \cdot 16x^4 \cdot 3y + 10 \cdot 8x^3 \cdot 9y^2 + 10 \cdot 4x^2 \cdot 27y^3 + 5 \cdot 2x \cdot 81y^4 + 1 \cdot 1 \cdot 243y^5\\
			= 32x^{5}+240x^{4}y+720x^{3}y^{2}+1080x^{2}y^{3}+810xy^{4}+243y^{5}$. 
		\end{enumerate}
	} 
\end{bt}

% nguồn: MaHoaNH_TuLuan_Lop10
\begin{bt}%[0D8B3-2]%[ma:46c65cdb]
	Tìm số hạng chứa $x^2$ trong khai triển  $(x+3)^5$.
	\loigiai{Số hạng cần tìm là $T_{k+1}=\mathrm{C}_5^kx^{5-k}\cdot 3^k$\\ 
	Theo bài ra $5-k=2$, vậy $k=3$ nên số hạng cần tìm là $10 \times  3^3\cdot x^2=270x^2$. }
\end{bt}

% nguồn: MaHoaNH_TuLuan_Lop10
\begin{bt}%[0D8K3-2]%[ma:bf81f027]
	Tìm hệ số chứa $ x^4 $ trong khai  triển $ (2x-1)^5 $.
	\loigiai{ 
	 Ta có
	\begin{eqnarray*}
		 (2x-1)^5 &=&\mathrm{C}_5^0(2x)^5-\mathrm{C}_5^1(2x)^4+\mathrm{C}_5^2(2x)^3-\mathrm{C}_5^3(2x)^2+\mathrm{C}_5^4(2x)-\mathrm{C}_5^5\\
		 &=& 32x^5-80x^4+80x^3-40x^2+10x-1.
	\end{eqnarray*}
	Dựa vào khai triển suy ra hệ số chứa $ x^4 $ là $ -80  $.
	} 
\end{bt}

% nguồn: MaHoaNH_TuLuan_Lop10
\begin{bt}%[0D8B3-2]%[ma:d2d0f57a]
	Sử dụng nhị thức Newton, khai triển biểu thức $(x^2-2)^5$ và tìm  hệ số của $x^4$ trong khi triển đó.
	 \loigiai{
		\allowdisplaybreaks 
		\begin{eqnarray*}
			 (x^2-2)^5&=&[x^2+(-2)]^5\\
			 &=&\mathrm{C}_5^0(x^2)^5+\mathrm{C}_5^1(x^2)^4(-2)+\mathrm{C}_5^2(x^2)^3(-2)^2+\mathrm{C}_5^3(x^2)^2(-2)^3+\mathrm{C}_5^4x^2(-2)^4+\mathrm{C}_5^5(-2)^5\\
			&=&x^{10}-10x^8+40x^6-80x^4+80x^2-32. 
		\end{eqnarray*}
		Do đó hệ số của $x^4$ trong khai  triển trên là $-80$.
	 }
\end{bt}
